% GENERATED FILE. Edit canonical lessons, then run npm run generate:books.
% canonical-source-hash: fnv1a64:91724edb5c781379
% canonical-lessons: GU-C242-pakvu, GU-C242-lanvu, GU-C242-cutvu, GU-C242-khodvu, GU-C242-gadhedo, GU-R242-review-237-242, GU-R242-review-final-242

\chapter{Harvest time}
\label{ch:gu-harvest-time}
\hlchaptermodality{242}

\begin{chapteropening}
\textbf{By the end of this chapter:} \emph{I can say what we harvested, picked and dug up.}

Complete the last lesson of chapter 242: I can say what we harvested, picked and dug up.
\end{chapteropening}

\section[\texorpdfstring{\gu{પાકવું} (pākvũ)}{પાકવું (pākvũ)}]{\gu{પાકવું} (pākvũ) — to ripen}

\label{lesson:GU-C242-pakvu}



\begin{quote}
Before the new one: say the Gujarati for to uproot, then the Gujarati for to die.
\end{quote}

\subsection*{You'll want to know: \gu{પાકવું}}
\textbf{\gu{પાકવું}} — \emph{pākvũ} — \textquotedblleft{}to ripen\textquotedblright{}.

\begin{quote}
\textbf{\gu{ખેતરમાં} \gu{ઘઉં} \gu{પાક્યા}.} — \emph{khetarmā̃ ghaũ pākyā} — \textquotedblleft{}The wheat ripened in the field.\textquotedblright{}
\end{quote}

The adjective is \textbf{\gu{પાકું}}, ripe.

\begin{cousinweb}
\textbf{\gu{પાકવું}} (\emph{pākvũ}) is a tadbhava. It is built on Sanskrit \emph{pakva}, \textquotedblleft{}cooked, ripe\textquotedblright{}, and CDIAL glosses the Gujarati verb \textquotedblleft{}to become ripe\textquotedblright{} (CDIAL 7621).
\end{cousinweb}

\begin{grammarlens}[title={What you've built}]
One more word for everyday life at A2.
\end{grammarlens}

\subsection*{Your turn}
\begin{itemize}
\raggedright
  \item \emph{Say it:} \emph{pākvũ}
  \item \emph{Say it:} \emph{pākvũ}, once more
  \item \emph{Say it:} \emph{khetarmā̃ ghaũ pākyā} — The wheat ripened in the field.
  \item \emph{Recall:} say the Gujarati for to uproot, then the Gujarati for to die, then say \emph{pākvũ} again
\end{itemize}

\begin{tcolorbox}[breakable,colback=teal!4,colframe=teal!35!black,title={Before you move on}]
What does \gu{પાકવું} mean? (\textquotedblleft{}To ripen\textquotedblright{}.) Say it once more.
\end{tcolorbox}
\section[\texorpdfstring{\gu{લણવું} (laṇvũ)}{લણવું (laṇvũ)}]{\gu{લણવું} (laṇvũ) — to harvest, to reap}

\label{lesson:GU-C242-lanvu}



\begin{quote}
Before the new one: say the Gujarati for to die, then the Gujarati for to ripen.
\end{quote}

\subsection*{You'll want to know: \gu{લણવું}}
\textbf{\gu{લણવું}} — \emph{laṇvũ} — \textquotedblleft{}to harvest, to reap\textquotedblright{}.

\begin{quote}
\textbf{\gu{અમે} \gu{ઘઉં} \gu{લણ્યા}.} — \emph{ame ghaũ laṇyā} — \textquotedblleft{}We harvested the wheat.\textquotedblright{}
\end{quote}

\begin{cousinweb}
\textbf{\gu{લણવું}} (\emph{laṇvũ}) is a tadbhava, a word inherited from Sanskrit and reshaped by sound change. It comes from \emph{lunāti}, \textquotedblleft{}cuts, reaps\textquotedblright{} (CDIAL 11082).
\end{cousinweb}

\begin{grammarlens}[title={What you've built}]
One more word for everyday life at A2.
\end{grammarlens}

\subsection*{Your turn}
\begin{itemize}
\raggedright
  \item \emph{Say it:} \emph{laṇvũ}
  \item \emph{Say it:} \emph{laṇvũ}, once more
  \item \emph{Say it:} \emph{ame ghaũ laṇyā} — We harvested the wheat.
  \item \emph{Recall:} say the Gujarati for to die, then the Gujarati for to ripen, then say \emph{laṇvũ} again
\end{itemize}

\begin{tcolorbox}[breakable,colback=teal!4,colframe=teal!35!black,title={Before you move on}]
What does \gu{લણવું} mean? (\textquotedblleft{}To harvest, to reap\textquotedblright{}.) Say it once more.
\end{tcolorbox}
\section[\texorpdfstring{\gu{ચૂંટવું} (cū̃ṭvũ)}{ચૂંટવું (cū̃ṭvũ)}]{\gu{ચૂંટવું} (cū̃ṭvũ) — to pick, to pluck}

\label{lesson:GU-C242-cutvu}



\begin{quote}
Before the new one: say the Gujarati for to ripen, then the Gujarati for to harvest.
\end{quote}

\subsection*{You'll want to know: \gu{ચૂંટવું}}
\textbf{\gu{ચૂંટવું}} — \emph{cū̃ṭvũ} — \textquotedblleft{}to pick, to pluck\textquotedblright{}.

\begin{quote}
\textbf{\gu{મેં} \gu{બગીચામાં} \gu{ફૂલ} \gu{ચૂંટ્યાં}.} — \emph{mẽ bagīcāmā̃ phūl cū̃ṭyā̃} — \textquotedblleft{}I picked flowers in the garden.\textquotedblright{}
\end{quote}

\begin{cousinweb}
\textbf{\gu{ચૂંટવું}} (\emph{cū̃ṭvũ}) is a tadbhava. CDIAL lists it, glossed \textquotedblleft{}to pluck out, select\textquotedblright{}, under \emph{cuṇṭati}, a by-form of Sanskrit \emph{cuṭati}, \textquotedblleft{}cuts off\textquotedblright{}, in a family it takes to be of non-Aryan origin (CDIAL 4857).
\end{cousinweb}

\begin{grammarlens}[title={What you've built}]
One more word for everyday life at A2.
\end{grammarlens}

\subsection*{Your turn}
\begin{itemize}
\raggedright
  \item \emph{Say it:} \emph{cū̃ṭvũ}
  \item \emph{Say it:} \emph{cū̃ṭvũ}, once more
  \item \emph{Say it:} \emph{mẽ bagīcāmā̃ phūl cū̃ṭyā̃} — I picked flowers in the garden.
  \item \emph{Recall:} say the Gujarati for to ripen, then the Gujarati for to harvest, then say \emph{cū̃ṭvũ} again
\end{itemize}

\begin{tcolorbox}[breakable,colback=teal!4,colframe=teal!35!black,title={Before you move on}]
What does \gu{ચૂંટવું} mean? (\textquotedblleft{}To pick, to pluck\textquotedblright{}.) Say it once more.
\end{tcolorbox}
\section[\texorpdfstring{\gu{ખોદવું} (khodvũ)}{ખોદવું (khodvũ)}]{\gu{ખોદવું} (khodvũ) — to dig}

\label{lesson:GU-C242-khodvu}



\begin{quote}
Before the new one: say the Gujarati for to harvest, then the Gujarati for to pick.
\end{quote}

\subsection*{You'll want to know: \gu{ખોદવું}}
\textbf{\gu{ખોદવું}} — \emph{khodvũ} — \textquotedblleft{}to dig\textquotedblright{}.

\begin{quote}
\textbf{\gu{મેં} \gu{બટાકા} \gu{ખોદ્યા}.} — \emph{mẽ baṭākā khodyā} — \textquotedblleft{}I dug up potatoes.\textquotedblright{}
\end{quote}

\begin{cousinweb}
\textbf{\gu{ખોદવું}} (\emph{khodvũ}) comes from \emph{*khodd}, \textquotedblleft{}dig\textquotedblright{}, a reconstructed form; CDIAL glosses the Gujarati verb \textquotedblleft{}to dig, carve\textquotedblright{} (CDIAL 3934).
\end{cousinweb}

\begin{grammarlens}[title={What you've built}]
One more word for everyday life at A2.
\end{grammarlens}

\subsection*{Your turn}
\begin{itemize}
\raggedright
  \item \emph{Say it:} \emph{khodvũ}
  \item \emph{Say it:} \emph{khodvũ}, once more
  \item \emph{Say it:} \emph{mẽ baṭākā khodyā} — I dug up potatoes.
  \item \emph{Recall:} say the Gujarati for to harvest, then the Gujarati for to pick, then say \emph{khodvũ} again
\end{itemize}

\begin{tcolorbox}[breakable,colback=teal!4,colframe=teal!35!black,title={Before you move on}]
What does \gu{ખોદવું} mean? (\textquotedblleft{}To dig\textquotedblright{}.) Say it once more.
\end{tcolorbox}
\section[\texorpdfstring{\gu{ગધેડો} (gadheḍo)}{ગધેડો (gadheḍo)}]{\gu{ગધેડો} (gadheḍo) — donkey}

\label{lesson:GU-C242-gadhedo}



\begin{quote}
Before the new one: say the Gujarati for to pick, then the Gujarati for to dig.
\end{quote}

\subsection*{You'll want to know: \gu{ગધેડો}}
\textbf{\gu{ગધેડો}} — \emph{gadheḍo} — \textquotedblleft{}donkey\textquotedblright{}.

\begin{quote}
\textbf{\gu{સાંજે} \gu{ગધેડો} \gu{ગામમાં} \gu{પાછો} \gu{આવ્યો}.} — \emph{sā̃je gadheḍo gāmmā̃ pācho āvyo} — \textquotedblleft{}In the evening the donkey came back to the village.\textquotedblright{}
\end{quote}

\begin{cousinweb}
\textbf{\gu{ગધેડો}} (\emph{gadheḍo}) is a tadbhava. It comes from Sanskrit \emph{gardabha}, \textquotedblleft{}ass\textquotedblright{}, where CDIAL lists Gujarati \emph{gaddho}, \textquotedblleft{}ass\textquotedblright{}; \emph{gadheḍo} adds the suffix \emph{-ḍo} (CDIAL 4054).
\end{cousinweb}

\begin{grammarlens}[title={What you've built}]
One more word for everyday life at A2.
\end{grammarlens}

\subsection*{Your turn}
\begin{itemize}
\raggedright
  \item \emph{Say it:} \emph{gadheḍo}
  \item \emph{Say it:} \emph{gadheḍo}, once more
  \item \emph{Say it:} \emph{sā̃je gadheḍo gāmmā̃ pācho āvyo} — In the evening the donkey came back to the village.
  \item \emph{Recall:} say the Gujarati for to pick, then the Gujarati for to dig, then say \emph{gadheḍo} again
\end{itemize}

\begin{tcolorbox}[breakable,colback=teal!4,colframe=teal!35!black,title={Before you move on}]
What does \gu{ગધેડો} mean? (\textquotedblleft{}Donkey\textquotedblright{}.) Say it once more.
\end{tcolorbox}
\section[(dialogue)]{Animals and plants - a review}

\label{lesson:GU-R242-review-237-242}



\begin{quote}
Say the Gujarati for crow and owl.
\end{quote}

\subsection*{Your turn}
Say these aloud:

\begin{itemize}
\raggedright
  \item crow, owl, wing, to flutter, to fly up
  \item cat, house lizard, spider, ant, to hang
  \item frog, mud, to sink, to sting, to slip away
  \item plant, seed, to plant, to irrigate, rose
  \item to bloom, to wilt, root, to uproot, to die
  \item to ripen, to harvest, to pick, to dig, donkey
\end{itemize}

\begin{tcolorbox}[breakable,colback=teal!4,colframe=teal!35!black,title={Before you move on}]
Crow? (\textbf{\gu{કાગડો}}.) Owl? (\textbf{\gu{ઘુવડ}}.) Wing? (\textbf{\gu{પાંખ}}.) To flutter? (\textbf{\gu{ફફડવું}}.) To fly up? (\textbf{\gu{ઊડવું}}.) Cat? (\textbf{\gu{બિલાડી}}.) House lizard? (\textbf{\gu{ગરોળી}}.) Spider? (\textbf{\gu{કરોળિયો}}.) Ant? (\textbf{\gu{કીડી}}.) To hang? (\textbf{\gu{લટકવું}}.) Frog? (\textbf{\gu{દેડકો}}.) Mud? (\textbf{\gu{કાદવ}}.) To sink? (\textbf{\gu{ડૂબવું}}.) To sting? (\textbf{\gu{ડંખવું}}.) To slip away? (\textbf{\gu{છટકવું}}.) Plant? (\textbf{\gu{છોડ}}.) Seed? (\textbf{\gu{બીજ}}.) To plant? (\textbf{\gu{રોપવું}}.) To irrigate? (\textbf{\gu{સીંચવું}}.) Rose? (\textbf{\gu{ગુલાબ}}.) To bloom? (\textbf{\gu{ખીલવું}}.) To wilt? (\textbf{\gu{કરમાવું}}.) Root? (\textbf{\gu{મૂળ}}.) To uproot? (\textbf{\gu{ઉખેડવું}}.) To die? (\textbf{\gu{મરવું}}.) To ripen? (\textbf{\gu{પાકવું}}.) To harvest? (\textbf{\gu{લણવું}}.) To pick? (\textbf{\gu{ચૂંટવું}}.) To dig? (\textbf{\gu{ખોદવું}}.) Donkey? (\textbf{\gu{ગધેડો}}.)
\end{tcolorbox}
\section[(dialogue)]{Harvest time — once more}

\label{lesson:GU-R242-review-final-242}



\begin{quote}
Say the Gujarati for to ripen and to harvest.
\end{quote}

\subsection*{Your turn}
\begin{itemize}
\raggedright
  \item \emph{Say it:} to ripen, to harvest, to pick, to dig, donkey
\end{itemize}

\begin{tcolorbox}[breakable,colback=teal!4,colframe=teal!35!black,title={Before you move on}]
To ripen? (\textbf{\gu{પાકવું}}.) To harvest? (\textbf{\gu{લણવું}}.) To pick? (\textbf{\gu{ચૂંટવું}}.) To dig? (\textbf{\gu{ખોદવું}}.) Donkey? (\textbf{\gu{ગધેડો}}.)
\end{tcolorbox}
